\documentclass[10pt,a4paper]{amsart}
\usepackage[margin=19mm]{geometry}
\usepackage{amsmath,amssymb,mathtools}
\usepackage[scr=rsfs,cal=euler]{mathalfa}
\usepackage{xfrac}
\usepackage[unicode,hidelinks,bookmarks=false]{hyperref}
\newcommand{\ie}{i.e.\ }
\newcommand{\cf}{cf.\ }
\newcommand{\resp}{respectively\ }
\newcommand{\Subsection}{\subsection}
\providecommand{\nobreakdash}{\nobreak\hbox{-}}
\setlength{\emergencystretch}{2em}
\multlinegap0pt
\usepackage{amssymb}
\newtheorem{lemma}{Lemma}
\newcommand{\T}{\mathcal{T}}
\newcommand{\treea}{{\rm tree}\textnormal{-}\alpha}
\title{On the Relation Between Treewidth, Tree-Independence Number, and Tree-Chromatic Number of Graphs}
\author{Alex Koutsoutis, Kilian Krause, Chun-Hung Liu, Mirza Redzic, Torsten Ueckerdt}
\date{Source: arXiv:2504.19751v2 (CC BY 4.0)}
\begin{document}
\maketitle

\noindent\emph{Source and licence.} On the Relation Between Treewidth, Tree-Independence Number, and Tree-Chromatic Number of Graphs. Version arXiv:2504.19751v2 (CC BY 4.0), distributed by arXiv under the Creative Commons Attribution 4.0 International licence, http://creativecommons.org/licenses/by/4.0/, as recorded in the arXiv licence metadata for this exact version and shown on the abstract page https://arxiv.org/abs/2504.19751v2. Full licence notice: \texttt{licenses/arxiv-2504.19751-CC-BY-4.0.txt}.\par\smallskip

\noindent\textit{Editorial scope.} Counted unit: the article's lemma lemma:subdivision-of-Kn-tree-alpha-lower-bound (source0.tex lines 412--414). The article's complete original proof of it is delivered with it (source0.tex lines 416--465, one \texttt{proof} environment of 50 lines). No mathematics is written, repaired or re-derived: the statement and the proof are the article's own lines, in the article's own notation, and no retained line is substituted or re-flowed.  No further source-local context is needed: every cross-reference of every retained line resolves inside the two delivered spans.  Nothing was omitted from any retained span, so the package carries no omission record.  No retained line contains a citation, so the wrapper carries no bibliography and no bibliography entry.  No retained line contains a figure environment, an image-inclusion command or drawing code, so no image is delivered and the package declares no asset dependency.  Editorial formatting only, in addition: an AMS \texttt{amsart} preamble that restates the article's own declarations and macros used by the retained lines, namely the environments \texttt{lemma} on separate counters, together with the standard packages those lines need.  The retained environments print in this document as Lemma 1 in order of appearance, whereas the article prints the same items with the numbering of its own full document.  The complete native source of the article as distributed in the arXiv e-print for this exact version is bundled unchanged as \texttt{sources/177194/source0.tex}.
\begin{lemma}\label{lemma:subdivision-of-Kn-tree-alpha-lower-bound}
    For every $n \geq 2$ we have $\treea(C(\overline{K_n})) = n-1$. %$\treea(S(K_n)) = n-1$.
\end{lemma}

\begin{proof}
    Let $A \subset V(C(\overline{K_n}))$ be the subset of $n$ vertices in $C(\overline{K_n})$ corresponding to the original vertices of $\overline {K_n}$.
    Note that $A$ is an independent set in $C(\overline{K_n})$.
    
    We first show the lower bound $\treea(C(\overline{K_n})) \geq n-1$.
    Consider a tree decomposition $\T = (T,\{X_t\}_{t\in V(T)})$ of $C(\overline{K_n})$ with $\alpha(\T) = \treea(C(\overline{K_n}))$.
    We may assume that no bag of $\T$ contains $A$, for otherwise $\alpha(\T) \geq |A| = n \geq n-1$, as desired.
    %If for some node $t \in V(T)$ it holds $A \subseteq X_t$, then $\alpha(\T) \geq |A| = n \geq n-1$, as desired.
    So there exists an edge $t_1t_2$ of $T$ such that $T - t_1t_2$ consists of two subtrees $T_1$ and $T_2$ of $T$, and there exists a partition $A_1 \cup A_2 \cup A_3$ of $A$ with the following properties:
    \begin{itemize}
        \item $|A_1| \geq 1$ and every vertex in $A_1$ appears only in bags of $T_1$.
        \item $|A_2| \geq 1$ and every vertex in $A_2$ appears only in bags of $T_2$.
        \item Every vertex in $A_3$ appears in bags of $T_1$ and bags of $T_2$.
    \end{itemize}
    Clearly, $|A_1|+|A_2|+|A_3| = |A| = n$.
    Moreover, the bag $X_{t_1}$ contains all vertices in $A_3$, as well as the set $B$ consisting of the (degree-$2$) vertices in $C(\overline{K_n})-A$ adjacent in $C(\overline{K_n})$ to both vertices in $A_1$ and $A_2$.
    Observe that $A_3 \cup B$ is an independent set in $C(\overline{K_n})$.
    So $\alpha(\T) \geq |A_3 \cup B| = |A_3| + |B|$.
    Since $|B| = |A_1|\cdot |A_2| \geq |A_1|+|A_2|-1$, where the inequalities uses that $|A_1|,|A_2| \geq 1$, we have $\alpha(\T) \geq |B| + |A_3| \geq |A_1|+|A_2|+|A_3|-1 = n-1$, as desired.

    \medskip

    Now we show the upper bound $\treea(C(\overline{K_n})) \leq n-1$.
    The case $n=2$ is trivial, so we may assume $n \geq 3$.
    Let $S$ be the star with $\binom{n-1}{2}+1$ leaves.
    We denote by $s$ the central vertex of $S$, and denote the leaves of $S$ by $s_0$ and $s_{i,j}$ for $1 \leq i <j \leq n-1$.
    Denote $A$ by $\{a_0,a_1,a_2,...,a_{n-1}\}$.
    Define $Y_s = (A-\{a_0\}) \cup N_{C(\overline{K_n})}(a_0)$ and $Y_{s_0} = \{a_0\} \cup N_{C(\overline{K_n})}(a_0)$.
    For every $i,j$ with $1 \leq i <j \leq n-1$, define $Y_{s_{i,j}}$ to be the set consisting of $a_i,a_j$ and the unique common neighbor of $a_i$ and $a_j$ in $C(\overline{K_n})$.
    Clearly, $(S,(Y_t)_{t \in V(S)})$ is a tree-decomposition of $C(\overline{K_n})$.
    Note that the subgraph induced by $Y_s$ contains a perfect matching with $n-1$ edges; the subgraph induced by each leaf of $S$ is a star with at most $\max\{|A|-1,2\} = n-1$ leaves.
    So every bag induces a subgraph with independence number at most $n-1$.
    This shows $\treea(C(\overline{K_n})) \leq n-1$.
%    , we construct a tree decomposition $\T = (T,\{X_t\}_{t\in V_T})$ of $S(K_n)$ with $\alpha(\T) = n-1$.
%    For this, pick one vertex $a \in A$ and let $F$ be the set of the $n-1$ edges in $K_n$ incident to $a$.
%    For $T$ we take a star with $\binom{n-1}{2}+1$ leaves and denote by $t$ the central vertex of $T$.
%    For the corresponding bag $X_t$, we take the union of $A - \{a\}$ and all subdivision vertices corresponding to edges in $F$.
%    Note that $S(K_n)[X_t]$ is a matching on $n-1$ edges.
%    
%    For each edge $e = uv$ in $E(K_n) - F$, let $t(e)$ be a corresponding leaf in $T$.
%    We define the corresponding bag as $X_{t(e)} = \{u,v,e\}$, where $e$ is the subdivision vertex of the edge $uv$ in $S(K_n)$.
%    Note that $S(K_n)[X_{t(e)}]$ is a path on two edges.
%    
%    Finally, we incorporate the special vertex $a \in A$ and its incident edges.
%    For this, let $t(a)$ be the last leaf of $T$ and define $X_{t(a)} = \{a\} \cup F$.
%    That is, $X_{t(a)}$ contains the original vertex $a$ of $K_n$ and the subdivision vertex of each edge $e$ incident to $a$ in $K_n$.
%    Note that $S(K_n)[X_{t(a)}]$ is a star with $n-1$ leaves.
%
%    Finally, observe that $\T = (T,\{X_t\}_{t\in V_T})$ is a valid tree decomposition of $S(K_n)$, and $\alpha(\T) = n-1$.
\end{proof}

\end{document}
